% Generated draft: requires mathematical and source review.
\documentclass[11pt]{article}
\usepackage[margin=25mm]{geometry}
\usepackage{fontspec}
\setmainfont{Latin Modern Roman}
\usepackage{amsmath,amssymb,mathtools,unicode-math}
\setmathfont{Latin Modern Math}
\usepackage{enumitem,needspace,xurl,hyperref}
\setlength{\parindent}{0pt}
\setlength{\parskip}{5pt}
\setlength{\emergencystretch}{3em}
\newcommand{\N}{\mathbb{N}}
\newcommand{\Z}{\mathbb{Z}}
\newcommand{\Q}{\mathbb{Q}}
\newcommand{\R}{\mathbb{R}}
\newcommand{\C}{\mathbb{C}}
\newcommand{\F}{\mathbb{F}}
\newcommand{\A}{\mathbb{A}}
\newcommand{\PP}{\mathbb{P}}
\newcommand{\E}{\mathbb{E}}
\newcommand{\eps}{\varepsilon}
\newcommand{\dd}{\,\mathrm{d}}
\begin{document}
\section*{UnsolvedMath ID: 11000310}
\textbf{Source problem code:} catalogue entry
\subsection{Definitions and mathematical statement}
Let \(F_n\) be free of rank \(n\) and \(\operatorname{Out}F_n=\operatorname{Aut}(F_n)/\operatorname{Inn}(F_n)\). The principal task is to construct nonzero rational homology or cohomology classes
\[
 x\in H_k(\operatorname{Out}F_n;\Q)
 \quad\text{or}\quad
 x\in H^k(\operatorname{Out}F_n;\Q),
\]
with ranks and positive degrees explicitly stated and nonvanishing justified. Positive degree is the substantive interpretation; degree-zero classes are automatic. The source does not require all ranks or a full ring computation.

The continuation defines the homologically trivial outer subgroup
\(I\operatorname{Out}_n=\ker(\operatorname{Out}F_n\to\operatorname{GL}(n,\Z))\).
It reports Igusa [38] constructing higher Franz--Reidemeister torsion classes
\[
 \tau_{2i}\in H^{4i}(I\operatorname{Out}_n;\R),
\]
and relates them to the source's Theorem 4.1: for the natural projection \(p\), the Borel classes satisfy
\[
 p^*\beta_{4i+1}=0\quad\text{in }H^{4i+1}(\operatorname{Out}F_n;\R),
 \qquad\beta_{4i+1}\in H^{4i+1}(\operatorname{GL}(n,\Z);\R).
\]
The exact rank/degree ranges and construction are omitted. The source says nontriviality of the \(\tau_{2i}\) was unknown; vanishing of the pulled-back Borel class alone is not asserted to prove that nontriviality.

The registry marks the item solved, but its summary says only partial classes are known and gives an unfinished \(H_2\) rank claim without an exact theorem. These reports do not identify a solution of the separate real-coefficient torsion question.
\subsection{Short English statement}
Construct positive-degree rational classes for outer automorphism groups, while keeping separate the reported real higher-torsion classes on their Torelli kernels and the question of their nonvanishing.
\subsection{Sources}
\begin{itemize}
\item \emph{Problems on Mapping Class Groups and Related Topics}, edited by Farb, Problem 4.3 and following discussion, supplied PDF page 365. \url{https://www.math.uchicago.edu/~farb/papers/mcgbook.pdf}. Recorded access 2026-07-29; extraction pdf-font-structure.
\item Igusa [38], higher torsion, and the preceding source Theorem 4.1 for Borel-class pullback vanishing. The bibliography, theorem hypotheses and class definitions are absent.
\item Local UnsolvedMath registry, \texttt{data/unsolvedmath/problems.json}, permanent ID \(11000310\), exact source code \texttt{catalogue entry}; solved/SOLVED-IN-LITERATURE label with a PARTIAL summary mentioning Morita/MMM-type work and an incomplete \(H_2\) assertion, without exact citations.
\end{itemize}
\end{document}
